\begin{mistake}{2026-05-28}{06}

\begin{problemcontent}
设积分区域
\[
D=\{(x,y)\mid |x|+|y|\le 1\},
\]
求二重积分
\[
\iint_D (2-x)(2-y)(1-|x|-|y|)\,d\sigma
\]
的值。选项：\((A)\ 1\)，\((B)\ \frac{4}{3}\)，\((C)\ 2\)，\((D)\ \frac{8}{3}\)。
\end{problemcontent}

\begin{solutioncontent}
记
\[
I=\iint_D (2-x)(2-y)(1-|x|-|y|)\,d\sigma .
\]
先展开
\[
(2-x)(2-y)=4-2x-2y+xy.
\]
因此
\[
\begin{aligned}
I={}&4\iint_D(1-|x|-|y|)\,d\sigma
-2\iint_Dx(1-|x|-|y|)\,d\sigma \\
&-2\iint_Dy(1-|x|-|y|)\,d\sigma
+\iint_Dxy(1-|x|-|y|)\,d\sigma .
\end{aligned}
\]
区域 \(D\) 关于 \(x\) 轴、\(y\) 轴均对称，且 \(1-|x|-|y|\) 关于 \(x,y\) 均为偶函数，所以
\[
\iint_Dx(1-|x|-|y|)\,d\sigma=0,
\quad
\iint_Dy(1-|x|-|y|)\,d\sigma=0,
\quad
\iint_Dxy(1-|x|-|y|)\,d\sigma=0.
\]
故
\[
I=4\iint_D(1-|x|-|y|)\,d\sigma.
\]
设
\[
J=\iint_D(1-|x|-|y|)\,d\sigma.
\]
由于被积函数关于 \(x,y\) 均为偶函数，可只算第一象限再乘 \(4\)。第一象限内 \(|x|=x, |y|=y\)，区域为
\[
0\le x\le 1,\quad 0\le y\le 1-x.
\]
于是
\[
\begin{aligned}
J&=4\int_0^1\int_0^{1-x}(1-x-y)\,dy\,dx \\
&=4\int_0^1\left[(1-x)y-\frac{y^2}{2}\right]_0^{1-x}\,dx \\
&=4\int_0^1\frac{(1-x)^2}{2}\,dx
=2\int_0^1(1-x)^2\,dx
=\frac{2}{3}.
\end{aligned}
\]
因此
\[
I=4J=4\cdot\frac{2}{3}=\frac{8}{3}.
\]
故答案为 \((D)\)。
\end{solutioncontent}

\mistakeknowledge{
\begin{itemize}
  \item \textbf{核心公式/定理}：若区域关于某轴对称，被积函数关于对应变量为奇函数，则该项二重积分为 \(0\)；若为偶函数，可取半区或四分之一区域再乘倍数。
  \item \textbf{易错点}：\(\sqrt{x^2}=|x|\)，不能误写为 \(x\)；本题不是化成 \(4xy\) 的积分，而是奇函数项消失后剩 \(4(1-|x|-|y|)\)。
  \item \textbf{关联知识}：\(|x|+|y|\le 1\) 是关于两坐标轴对称的菱形区域；第一象限内化为 \(x\ge0,\ y\ge0,\ x+y\le1\)。
\end{itemize}}

\end{mistake}
